% Zinsrechnung – bank/zinsrechnung/ (zusammenbau v0.4, Heft)
\einheitenkopf[t4]{Zinsrechnung}
\setcounter{aufgabe}{14}

\begin{aufgabe}{Zinsen berechnen – Prüfungsaufgaben}
\begin{teile}
\teil Ein Sparbetrag von 2\,600\,€ wird ein Jahr lang mit 2\,\% pro Jahr verzinst. Wie viel Euro Zinsen gibt es nach einem Jahr? \hfill (P10 2015 OS) \\ \leerfeld[€]
\teil Eine Geldanlage von 4\,500\,€ wird ein Jahr lang mit 3\,\% pro Jahr verzinst. Wie viel Euro Zinsen gibt es nach einem Jahr? \hfill (P10 2015 OS) \\ \leerfeld[€]
\teil Ein Sparguthaben von 15\,000\,€ bringt in einem Jahr 270\,€ Zinsen. Gib den Zinssatz an. \hfill (P10 2014 OS) \\ \leerfeld[\%]
\teil Ein Sparguthaben von 2\,500\,€ bringt in einem Jahr 65\,€ Zinsen. Gib den Zinssatz an. \hfill (P10 2014 OS) \\ \leerfeld[\%]
\end{teile}
\end{aufgabe}
